% ======================================================================
% SECTION 2: METHODS
% Six thematic subsections: instruction inputs, robot 2030 NeuroSpeed,
% sensors and sample table, xyz coordinate mapping, iterations plus
% LLM tournament, code execution environment.
% Location: 2030-gbm-1min/paper/full-paper/sections/methods.tex
% ======================================================================

\subsection{Instruction Inputs}
\label{sec:methods-inputs}

The 2030-gbm-1min repository was generated from 12 hand-authored
instruction files at
kevinkawchak/physical-ai-oncology-trials/competitions/instructions/one_minute_variant/
\cite{one-minute-variant-instructions}. Together the 12 files amount
to roughly 130 KB of read-only context; a senior author can read and
verify all 12 in approximately one hour. The complete inventory is:
README.md (scope and per-commit roadmap), commit_01_overview_1min.md
(project narrative and per-commit roadmap),
commit_02_sensors_1min.md (sensor specification at 200 channels
across 4 arms mixed at 1 kHz commands plus 10 kHz force),
commit_03_xyz_4arm.md (deterministic sensor to xyz transformation
across 4 cooperating arms), commit_04_iterations_1min.md (16-iteration
deterministic sweep design and YAML configuration),
commit_05_competition_1min.md (LLM tournament and composite score
formula), file_size_pyramid_1min.md (Git commit caps and Zenodo
deposition strategy), glioblastoma_context_1min.md (IDH-wildtype
glioblastoma clinical context), multi_arm_coordination.md (1 kHz
heartbeat broadcast bus protocol), robot_specification_neurospeed.md
(2030 Medtronic NeuroSpeed 1.0 7-DOF DH parameters and per-arm tool
assignment), sensor_specification_10khz.md (50 channels per arm at
1 kHz commands plus 10 kHz force), and zenodo_archive_protocol.md
(Zenodo L0 raw deposition protocol at DOI
10.5281/zenodo.18445179). The entire kevinkawchak/robotic-surgeries
2030-gbm-1min/ tree was authored by Claude Code Opus 4.7 1M Max
\cite{claude-opus-47, claude-code} in a single sequence of seven
commits in one pull request. Approximations made during the
generation pass (count of files, file sizes, sampling rate
selections, iteration count) are accounted for in
\S\ref{sec:limitations}, particularly
\S\ref{sec:limits-accounting}.

\subsection{2030 Medtronic NeuroSpeed 1.0 Robot}
\label{sec:methods-robot}

The hypothetical 2030 Medtronic NeuroSpeed 1.0 multi-arm parallel
stereotactic neurosurgical robot is composed of 4 cooperating arms
with 7 degrees of freedom each (28 DOF total), a 0.5 m radius
hemisphere workspace shared across the 4 arms, and a hybrid
ultrasonic plus waterjet plus pulsed plasma tissue removal mechanism
on arm 1 at 800 mm cubed per second peak. The full per-arm 7-DOF
Denavit-Hartenberg parameter table reads as in
Table \ref{tab:dh}.

\begin{table}[H]
\centering
\caption{Per-arm 7-DOF Denavit-Hartenberg parameters for the 2030
Medtronic NeuroSpeed 1.0 (modified DH convention; identical across
all 4 arms). Joint 1 is the shoulder base; joint 7 is the tool flange.}
\label{tab:dh}
\begin{tabular}{>{\raggedright\arraybackslash}p{1.5cm}
                >{\raggedright\arraybackslash}p{2.0cm}
                >{\raggedright\arraybackslash}p{2.0cm}
                >{\raggedright\arraybackslash}p{2.0cm}
                >{\raggedright\arraybackslash}p{2.0cm}
                >{\raggedright\arraybackslash}p{2.5cm}}
\toprule
\textbf{Joint i} & \textbf{a (mm)} & \textbf{alpha (rad)} & \textbf{d (mm)} & \textbf{theta range (rad)} & \textbf{vel max (rad/s)} \\
\midrule
1 & 0   & pi/2  & 340  & [-pi, pi]         & 6.28 \\
2 & 0   & -pi/2 & 0    & [-pi/2, pi/2]     & 6.28 \\
3 & 0   & pi/2  & 400  & [-pi, pi]         & 6.28 \\
4 & 0   & -pi/2 & 0    & [-pi/2, pi/2]     & 6.28 \\
5 & 0   & pi/2  & 400  & [-pi, pi]         & 6.28 \\
6 & 0   & -pi/2 & 0    & [-pi/2, pi/2]     & 6.28 \\
7 & 0   & 0     & 126  & [-2pi, 2pi]       & 6.28 \\
\bottomrule
\end{tabular}
\end{table}

Per-arm tip force is capped at 5.0 N and lateral force at 1.0 N
\cite{iec-80601-2-77}. The cumulative 4-arm tip force on the patient
frame is capped at 12.0 N. The E-stop budget is 5 ms; the emergency
arm-park trigger latency is 100 microseconds; the inter-arm minimum
clearance is 8 mm. The 1 kHz heartbeat broadcast bus carries a 32-
byte status frame per arm with a 1 ms deadline and a 3 ms watchdog
\cite{iec-62304}. The per-arm tool assignment is fixed: arm 1 hybrid
ultrasonic plus waterjet plus pulsed plasma, arm 2 bipolar coagulation
plus irrigation, arm 3 suction plus tissue collection, arm 4 0.5 T
intraoperative magnetic resonance imaging (iMRI) plus 5-ALA
fluorescence plus ultrasound. The 4-phase 60-second timeline is:
Phase 1 dural opening final 0 to 5 s, Phase 2 bulk resection 5 to
45 s at 800 mm cubed per second peak, Phase 3 margin and fine
resection 45 to 55 s, Phase 4 hemostasis verification and arm
withdrawal 55 to 60 s. Pre-operative anesthesia, registration, dural
opening, and multi-arm setup are precomputed and frozen at T-1800 s
to T+0 s relative to the 60-second window. The baseline ROSA ONE
Brain v3.0 is referenced for comparison \cite{rosa-one-brain}.

\subsection{Sensor Specification (54 columns by 1001 rows)}
\label{sec:methods-sensors}

The per-arm sensor schema carries 50 channels at mixed sample rates:
21 joint kinematics channels at 1 kHz (7 position plus 7 velocity
plus 7 torque), 7 end-effector pose channels at 1 kHz (3 position
plus 4 scalar-first quaternion), 6 end-effector force and torque
channels at 10 kHz (3 force plus 3 torque), 3 navigation deviation
channels at 1 kHz, 7 tool flag and adjunct channels at 1 kHz, and 6
safety enumeration and metadata channels at 1 kHz. Across 4
cooperating arms the total channel count is 200. The JSON Schema,
Protocol Buffers proto, and Apache Avro avsc declarations live at
kevinkawchak/robotic-surgeries/2030-gbm-1min/schemas/sensor_record_4arm.\{schema.json, proto, avsc\}.
The release-aggregate L0 raw archive (416 MB across 16 iterations)
is deposited on Zenodo only at DOI 10.5281/zenodo.18445179
\cite{zenodo, repo-robotic-surgeries}; the committed repository
carries a 1 KB Zenodo pointer JSON at
2030-gbm-1min/data/sensor_l0_raw_4arm.zenodo_pointer.json. The L1
through L3 aggregates use Apache Parquet with zstd-3 compression
\cite{apache-arrow}.

The 54 column by 1001 row sensor sample table at
kevinkawchak/robotic-surgeries/2030-gbm-1min/outputs/sensors/sensor_sample_4arm.csv
is the exceptional processing feat that no human team could produce
under the time budget of this paper. The 54 columns hold 50 sensor
channels per arm collapsed plus 4 metadata columns (timestamp_us,
arm_id, phase, command_state). The 1001 rows decimate the 10 kHz
force stream to a 1 kHz alignment with the command tick. A
representative slice is shown in Table \ref{tab:sensor-slice}; the
full file is 328,639 bytes; the JSONL companion at
outputs/sensors/sensor_sample_4arm.jsonl is 891,992 bytes and
carries the full per-arm record without column collapse.

\begin{table}[H]
\centering
\caption{Representative slice of the 54 column by 1001 row sensor
sample table; 8 of 54 columns and 6 of 1001 rows from Phase 2 bulk
resection on arms 1 and 4.}
\label{tab:sensor-slice}
\begin{tabular}{>{\raggedright\arraybackslash}p{1.5cm}
                >{\raggedright\arraybackslash}p{0.9cm}
                >{\raggedright\arraybackslash}p{0.9cm}
                >{\raggedright\arraybackslash}p{1.4cm}
                >{\raggedright\arraybackslash}p{1.4cm}
                >{\raggedright\arraybackslash}p{1.4cm}
                >{\raggedright\arraybackslash}p{1.5cm}
                >{\raggedright\arraybackslash}p{1.4cm}}
\toprule
\textbf{ts (us)} & \textbf{arm} & \textbf{phase} & \textbf{ee_x (mm)} & \textbf{ee_y (mm)} & \textbf{ee_z (mm)} & \textbf{tip force (N)} & \textbf{state} \\
\midrule
12000000 & 1 & 2 & 142.31 & -8.07 & 102.55 & 3.81 & EMIT \\
12000000 & 4 & 2 & 105.66 & 22.10 & 158.71 & 0.04 & EMIT \\
12001000 & 1 & 2 & 142.95 & -8.10 & 102.40 & 3.92 & EMIT \\
12001000 & 4 & 2 & 105.67 & 22.11 & 158.70 & 0.04 & EMIT \\
12002000 & 1 & 2 & 143.59 & -8.14 & 102.25 & 4.05 & EMIT \\
12002000 & 4 & 2 & 105.69 & 22.13 & 158.68 & 0.05 & EMIT \\
\bottomrule
\end{tabular}
\end{table}

\subsection{Per-Arm xyz Coordinate Mapping}
\label{sec:methods-xyz}

The deterministic sensor to xyz transformation is implemented in
kevinkawchak/robotic-surgeries/2030-gbm-1min/src/mapping/sensor_to_xyz_4arm.py
(Python reference) and in
2030-gbm-1min/src/control/robot_loop_4arm.cpp (C++20 real-time
control loop with a 1 ms command-to-actuator budget). Forward
kinematics maps the 7-DOF joint state to an end-effector pose in
the patient frame; inverse kinematics resolves an xyz command back
into a joint-target trajectory at a configurable tolerance in the
set \{1e-6, 1e-5, 1e-4, 1e-3\}. The 1 kHz heartbeat coordination
layer is implemented in 2030-gbm-1min/src/coordination/arm_heartbeat.cpp
with 32-byte status frames, a 3 ms watchdog, and the cumulative
12 N tip-force cap enforced at the heartbeat boundary
\cite{iec-80601-2-77, iec-62304}. The xyz command schema is at
2030-gbm-1min/schemas/xyz_command_4arm.schema.json.

The per-arm xyz command sample is split into 4 files at
2030-gbm-1min/outputs/xyz_mapping/xyz_trace_sample_arm1.csv through
xyz_trace_sample_arm4.csv (60 rows per arm, 240 commands total) with
an ASCII per-second xyz path overlay at
outputs/xyz_mapping/xyz_path_4arm.txt. Table \ref{tab:cmd-state} gives
the per-arm command-state distribution. All 240 commands resolve to
command_state=EMIT in the sample window; no rejections and no E-stop
triggers are observed.

\begin{table}[H]
\centering
\caption{Per-arm command-state distribution across the 240-row xyz
command sample at outputs/xyz_mapping/. All commands resolve to EMIT.}
\label{tab:cmd-state}
\begin{tabular}{>{\raggedright\arraybackslash}p{1.6cm}
                >{\raggedright\arraybackslash}p{1.8cm}
                >{\raggedright\arraybackslash}p{1.4cm}
                >{\raggedright\arraybackslash}p{1.4cm}
                >{\raggedright\arraybackslash}p{1.4cm}
                >{\raggedright\arraybackslash}p{1.4cm}
                >{\raggedright\arraybackslash}p{1.6cm}}
\toprule
\textbf{Arm} & \textbf{Tool} & \textbf{EMIT} & \textbf{HOLD} & \textbf{REJECT} & \textbf{ESTOP} & \textbf{Phase mix} \\
\midrule
1 & hyb u-w-p   & 60 & 0 & 0 & 0 & P1-P4 \\
2 & bipolar    & 60 & 0 & 0 & 0 & P1-P4 \\
3 & suction    & 60 & 0 & 0 & 0 & P2-P4 \\
4 & iMRI/5-ALA & 60 & 0 & 0 & 0 & P1-P4 \\
\bottomrule
\end{tabular}
\end{table}

\subsection{16-Iteration Sweep and On-Premises LLM Tournament}
\label{sec:methods-iterations}

The 16-iteration deterministic sweep is configured in
kevinkawchak/robotic-surgeries/2030-gbm-1min/config/iterations.yaml and
executed by 2030-gbm-1min/src/simulation/iterate_1min.py
(orchestrator, Python with concurrent.futures.ProcessPoolExecutor)
and 2030-gbm-1min/src/simulation/runner_1min.rs (Rust 2021 inner
loop with rand_pcg deterministic PRNG and crossbeam-channel
inter-arm coordination). The sweep covers four dimensions: noise
(0.01 to 0.05 mm sensor noise sigma, linear), force feedback gain
(0.8 to 1.2, linear), inverse-kinematics tolerance (1e-6 to 1e-3,
logarithmic), and heartbeat jitter (0 to 50 microseconds, linear).
The fixed seed is 20260510. The pipeline is bit-deterministic for
fixed seed plus parameter tuple.

The composite score formula is frozen at quality 0.40, time 0.25,
cost 0.20, safety 0.10, patient experience 0.05; the per-iteration
metric records are computed by
2030-gbm-1min/src/metrics/compute_1min.py. The on-premises LLM
tournament is implemented at 2030-gbm-1min/src/llm/compare_agent_1min.py
with a versioned prompt at 2030-gbm-1min/prompts/comparison_prompt_1min.md.
The default backend is Anthropic Claude Opus 4.7
\cite{claude-opus-47, claude-code}; on-premises Ollama \cite{ollama}
and vLLM \cite{vllm} are supported as drop-in alternatives. The
tournament runs a 4-entity bracket in two configurations: the default
robot-vs-robot bracket at 2030-gbm-1min/outputs/comparison/ (6
rounds) and the mixed 2 robot plus 2 human bracket at
2030-gbm-1min/outputs/comparison_robot_vs_human/ (6 rounds). The
structural-time-dimension caveat (1-minute robot vs 1-hour human
baseline) is preserved in every LLM judge rationale.

Per-iteration outputs land in
kevinkawchak/robotic-surgeries/2030-gbm-1min/outputs/iterations/ and
consist of an L1 50 ms aggregate Parquet, an L2 1 s aggregate
Parquet, an L3 per-phase aggregate Parquet, an events Parquet, and
a 1 KB L0 Zenodo pointer JSON per iteration. The cross-iteration
index.jsonl manifest plus the aggregate.duckdb analytical store
\cite{duckdb} support cross-iteration queries.

\subsection{Code Execution Environment}
\label{sec:methods-execution}

The repository is exercised on Linux (Ubuntu 22.04 LTS or later),
MacOS (Apple Silicon M3 Ultra or Intel), and Windows 11 with
PowerShell 7. An optional NVIDIA A100 GPU recipe (cargo build
--release --features cuda) accelerates the Rust runner. Claude
Code (CLI, web, and IDE plugin) executes the end-to-end pipeline
without any local Python or Rust install beyond what the recipe
requires \cite{claude-code, claude-opus-47}. Table \ref{tab:platforms}
summarizes the expected single-iteration wall-clock per platform.

\begin{table}[H]
\centering
\caption{Single-iteration wall-clock per platform under the
deterministic seed 20260510. M3 Ultra and A100 figures are read
from 2030-gbm-1min/README.md.}
\label{tab:platforms}
\begin{tabular}{>{\raggedright\arraybackslash}p{4.0cm}
                >{\raggedright\arraybackslash}p{3.8cm}
                >{\raggedright\arraybackslash}p{2.6cm}
                >{\raggedright\arraybackslash}p{2.6cm}}
\toprule
\textbf{Platform} & \textbf{Toolchain} & \textbf{Wall-clock} & \textbf{Notes} \\
\midrule
Linux Ubuntu 22.04 LTS & Python 3.10/3.11/3.12, Rust 2021, g++ & 60 s & x86_64 32-core reference \\
MacOS Apple Silicon M3 Ultra & Python 3.12, Rust 2021, clang++ & 26 to 32 s & faster than reference \\
Windows 11 PowerShell 7 & Python 3.11, Rust 2021, MSVC & 60 to 75 s & WSL2 optional \\
NVIDIA A100 GPU & Rust + cargo --features cuda & 12 s & cuda 12.x; optional \\
Claude Code (CLI/web/IDE) & managed runtime & varies & cloud or local backend \\
\bottomrule
\end{tabular}
\end{table}

CI gates run ruff format --check, ruff check, and yamllint -d
relaxed on Python 3.10, 3.11, and 3.12. The C++20 control loop and
C++20 1 kHz heartbeat layer compile on Linux (g++), MacOS (clang++),
and Windows (MSVC); the Rust 2021 runner builds under stable rustc.
The verification recipe in 2030-gbm-1min/README.md walks the user
through running iterate_1min.py for a single iteration, inspecting
the L1 to L3 Parquet output, and running compare_agent_1min.py with
the Ollama backend if the Anthropic API is not configured.
